\chapter{Implementation and Results}
\label{chap:results}

\noindent
This chapter details the software and hardware environments used,
presents the full experimental results across all twenty-three
configurations, analyses the findings, and discusses their implications
for probabilistic electricity demand forecasting.
All experiments share an identical data pipeline, evaluation protocol,
and random-seed management scheme, ensuring that reported differences
reflect genuine model-level contributions.

% ===========================================================================
\section{Environment and Working Tools}
\label{sec:environment}
% ===========================================================================

\subsection{Hardware Environment}
\label{subsec:hardware}

Three computing environments were used throughout the project
(Table~\ref{tab:hardware}).
Machine~1 (RTX~3070, 8~GB VRAM) was the primary development and
prototyping environment; mixed-precision \texttt{fp16} training was
necessary to fit the exogenous feature models within the 8~GB budget.
Machine~2 (RTX~4070~Ti, 12~GB VRAM) provided intermediate-scale validation
of key configurations before committing to multi-seed production runs.
The CCK virtual server (NVIDIA A40, 24~GB VRAM) was the sole environment
for all production multi-seed experiments, ensuring numerical precision
and full reproducibility.

\begin{table}[H]
\centering
\caption{Hardware environments used in this project.}
\label{tab:hardware}
\begin{tabular}{llllll}
\toprule
\textbf{Unit} & \textbf{GPU} & \textbf{VRAM} & \textbf{RAM} &
\textbf{Precision} & \textbf{Role} \\
\midrule
Work Unit 1 & RTX 3070    &  8 GB & 32 GB & fp16 & Dev / prototyping \\
Work Unit 2 & RTX 4070 Ti & 12 GB & 32 GB & fp32 & Intermediate validation \\
CCK Server  & NVIDIA A40  & 24 GB & 64 GB & fp32 & Production multi-seed runs \\
\bottomrule
\end{tabular}
\end{table}

\textbf{Training time summary:}
\begin{itemize}
  \item Baseline NBEATSAQCAT (A40, fp32, batch~2048): $\approx$8~hours per
        full 15-epoch run per seed.
  \item Exogenous variants (A40, fp32, batch~1024): $\approx$14~hours per
        seed due to the larger VRAM footprint of the feature encoders.
  \item Exogenous variants (RTX~3070, fp16, batch~512): 20--24~hours per seed.
  \item Novel architectural head variants (A40, fp32): 8--10~hours per seed.
\end{itemize}

% ---------------------------------------------------------------------------
\subsection{Software Environment}
\label{subsec:software}

\subsubsection{Programming Languages}
\label{subsubsec:languages}

Python~3.10 was the sole language for all machine-learning and
data-processing components.  Its ecosystem---NumPy, SciPy, Pandas,
Matplotlib, and Plotly---provides the numerical and visualisation
primitives required for large-scale time-series experimentation.
TypeScript~5 was used for the entire frontend web application, providing
compile-time type safety over the React component tree and API client code.

\subsubsection{Frameworks and Libraries}
\label{subsubsec:frameworks}

\textbf{Deep Learning:}
\begin{itemize}
  \item \textbf{PyTorch~2.0+}: automatic differentiation, CUDA kernel
        dispatch, and mixed-precision tensor operations.
        \texttt{torch.compile} was evaluated but disabled in production
        due to incompatibilities with PyTorch Lightning callbacks.
  \item \textbf{PyTorch Lightning~1.8+}: training loop orchestration,
        gradient accumulation, AMP, and the \texttt{ModelCheckpoint}
        callback for best-epoch model selection across all seeds.
  \item \textbf{OmegaConf / Hydra}: type-safe YAML configuration parsing
        enabling hierarchical experiment configuration with CLI overrides.
\end{itemize}

\textbf{Data and Evaluation:}
\begin{itemize}
  \item \textbf{NumPy / SciPy / Pandas}: data loading, pre-processing,
        and metric computation (CRPS, coverage, MAPE, sMAPE).
  \item \textbf{GluonTS}: time-series batching utilities and baseline
        implementations (DeepAR, WaveNet, TFT, Transformer, MLP).
  \item \textbf{scikit-learn}: Quantile Regression baseline
        (\texttt{QuantileRegressor}) and preprocessing utilities.
\end{itemize}

\textbf{Web Application:}
\begin{itemize}
  \item \textbf{React~18 + TypeScript}: component-based, type-safe UI with
        Recharts for interactive probabilistic forecast visualisations.
  \item \textbf{Vite}: fast development server and production bundler.
  \item \textbf{TailwindCSS}: responsive, utility-first styling.
  \item \textbf{Python FastAPI}: asynchronous REST back-end serving of
        pre-loaded AQNBEATSPlusPlus checkpoint; a single forward pass handles
        arbitrary quantile-level sets at inference time.
\end{itemize}

\subsubsection{Database}
\label{subsubsec:database}

\begin{itemize}
  \item \textbf{PostgreSQL~15+}: primary relational store for time-series
        data, user records, forecast results, and audit logs.
        All tables are indexed on \texttt{(country, timestamp)} for
        efficient range queries over historical demand.
  \item \textbf{Redis}: in-memory cache for real-time dashboard state
        and asynchronous training job queuing via Celery workers.
  \item \textbf{CSV experiment logs}: all results stored in structured
        CSV files (\texttt{exp\_name, seed, epoch, crps, ncrps, coverage,
        mape, mae, smape}) for post-hoc analysis and reproducibility audits.
\end{itemize}

\subsubsection{Development Tools}
\label{subsubsec:devtools}

\begin{itemize}
  \item \textbf{Git / GitHub}: version control using a
        branch-per-experiment strategy; each configuration change is
        tracked and tagged with its best-seed CRPS.
  \item \textbf{Docker / Docker Compose}: containerised environment
        eliminating dependency conflicts across the three machines.
  \item \textbf{TensorBoard}: real-time training-curve visualisation
        (loss, learning rate, validation CRPS and coverage by country).
  \item \textbf{Jupyter Notebooks}: interactive result analysis,
        prediction-interval plots, and ablation study visualisations.
  \item \textbf{Postman}: REST API testing and endpoint documentation.
  \item \textbf{VS~Code}: primary IDE with Python, Pylance, ESLint,
        and YAML extensions.
\end{itemize}

% ===========================================================================
\section{Computational Results}
\label{sec:comp_results}
% ===========================================================================

All configurations are evaluated on the EMHIRES 2018 held-out test set
(January--December~2018) unless otherwise stated.
The primary metric is the Continuous Ranked Probability Score
(CRPS,~Eq.~\ref{eq:crps_quantile}); lower is better.
Coverage at $\alpha{=}0.95$ (Eq.~\ref{eq:coverage_formal}) is reported
alongside CRPS, as calibration quality is equally important to sharpness
in operational forecasting.

% ---------------------------------------------------------------------------
\subsection{Baseline Reproduction and Exogenous Feature Augmentation}
\label{subsec:baseline_exog}

\subsubsection{Baseline Reproduction}

Experiment~3 reproduces the AQ-NBEATS baseline
\cite{smyl2025anyquantile} using the published hyperparameters.
With 7 random seeds on the CCK A40, we obtain
$\mathrm{CRPS} = 190.78 \pm 1.42$, $\mathrm{MAPE} = 2.376\%$, and
coverage-0.95 $= 0.942 \pm 0.011$.
The CRPS is approximately $9.6\%$ below the published 211.22, consistent
seeds.  The mild under-coverage (0.942 versus the 0.95 target) motivates
the CQR post-processing step in Section~\ref{sec:cqr_results}.

\subsubsection{Exogenous Feature Augmentation}

The pinball training objective is minimised
over the multi-quantile loss $\mathcal{L}_{\mathrm{MQ}}$
on input tuples that include the exogenous feature
vector, concatenated to the historical load
window via the feature fusion layer.

\textbf{Exogenous features alone (Exp~20)} achieves
$\mathrm{CRPS} = \mathbf{168.46}$, a $\approx$20.3\% reduction relative
to the published baseline.  This is the single largest improvement found in
this study.  Three training adjustments were necessary for convergence:
(i)~learning rate reduced from $5 \times 10^{-4}$ to $1 \times 10^{-4}$;
(ii)~batch size reduced from 1024 to 512 (VRAM constraint, fp16);
(iii)~early stopping at epoch~12 of~15.

\textbf{Series embeddings only (Exp~19)} achieves
$\mathrm{CRPS} = 206.51 \pm 0.88$ across 7 seeds, a modest 2.2\%
improvement.  The 32-dimensional per-country embeddings capture
demand-level biases but provide no direct access to meteorological drivers.

\textbf{Exogenous + series embeddings (Exp~21)} achieves
$\mathrm{CRPS} = 174.09 \pm 0.95$ with 5 seeds.  The embedding scale
parameter was reduced from 0.15 to 0.08 to prevent the low-dimensional
country embedding from competing with the richer exogenous signal.

\textbf{Combined with monotonicity penalty (Exp~22)} failed to improve
upon Exp~21: $\mathrm{CRPS} = 239.22$, coverage $= 0.985$.
The soft non-crossing penalty overwhelmed the exogenous training signal,
forcing excessively wide intervals.  This constitutes an important negative
finding for future practitioners.

\begin{table}[H]
\centering
\caption{Exogenous feature and series embedding results.
  CRPS lower is better; nominal coverage target is 0.95.}
\label{tab:exog_results}
\begin{tabular}{llccccc}
\toprule
\textbf{Exp} & \textbf{Configuration} & \textbf{Seeds} &
\textbf{CRPS} & \textbf{Cov-0.95} & \textbf{MAPE (\%)} &
\textbf{$\Delta$CRPS} \\
\midrule
Paper  & AQNBEATSPlusPlus (published)     & 8 & 211.22            & ---             & 2.47 & ---        \\
Exp~3  & Baseline reproduced       & 7 & $190.78\pm1.42$   & $0.942\pm0.011$ & 2.38 & $-9.6\%$   \\
Exp~19 & Series embed only         & 7 & $206.51\pm0.88$   & $0.936\pm0.002$ & 2.44 & $-2.2\%$   \\
Exp~20 & Exog only (1 seed)        & 1 & \textbf{168.46}   & 0.919           & 2.11 & $-20.3\%^*$\\
Exp~21 & Exog + Series             & 5 & $174.09\pm0.95$   & $0.912\pm0.009$ & 2.13 & $-17.6\%$  \\
Exp~22 & Exog+Series+Monotone      & 5 & 239.22            & 0.985           & 2.89 & $+13.2\%$  \\
\bottomrule
\multicolumn{7}{l}{\footnotesize $^*$Single seed; multi-seed mean may differ.}
\end{tabular}
\end{table}

% ---------------------------------------------------------------------------
\subsection{Curriculum Learning and Quantile Sampling}
\label{subsec:curriculum_results}

\subsubsection{Adaptive Sampling Without Curriculum}

Experiments~23--25 applied adaptive quantile sampling
directly on the Exp~21 configuration,
without any curriculum stage.  All three variants degraded performance
(Table~\ref{tab:adaptive}), with coverage dropping to 0.866--0.874.
The negative result establishes a key dependency: adaptive sampling is
only beneficial when the model already exhibits reasonable tail coverage.
Without it, the scheme over-concentrates on erratically shifting
calibration-gap estimates during early training, inducing gradient
instability.

\begin{table}[H]
\centering
\caption{Effect of adaptive quantile sampling applied directly on Exp~21
  (all 5 seeds).}
\label{tab:adaptive}
\begin{tabular}{lccc}
\toprule
\textbf{Configuration} & \textbf{CRPS} & \textbf{Cov-0.95} & \textbf{Status} \\
\midrule
Exog+Series (Exp~21) reference           & $174.09\pm0.95$ & $0.912\pm0.009$ & Reference \\
+Adaptive ($5q$, $\gamma{=}0.8$)         & 177.60          & 0.866           & Degraded  \\
+Adaptive ($3q$, $\gamma{=}1.0$)         & 178.03          & 0.874           & Degraded  \\
+Adaptive ($2q$, $\gamma{=}1.2$, s=0.05) & 177.53         & 0.868           & Degraded  \\
\bottomrule
\end{tabular}
\end{table}

\subsubsection{Curriculum v1: Feature-Staging}

Curriculum~v1 builds the model progressively across three stages, each
resuming from the previous checkpoint.
Stage~1 (exogenous only) achieves CRPS~173.0; Stage~2 adds series
embeddings, paradoxically increasing CRPS to~181.0; Stage~3 adds
adaptive sampling, further degrading to~190.0 with poor coverage (0.885).
The Stage~2 deterioration indicates that introducing embeddings on top of
a trained exogenous encoder requires a warm-up period that the fixed-epoch
schedule does not provide.

\subsubsection{Curriculum v2: Sampling-First (Best Configuration)}
\label{subsubsec:curriculum_v2}

Curriculum~v2 reverses the priority, establishing tail calibration
\emph{before} sharpening the distribution:
\begin{itemize}
  \item \textbf{Stage~1} (Exog + Beta sampling,
        Eq.~\ref{eq:beta_sampling}): quantile levels sampled from
        $\tau \sim \mathrm{Beta}(\beta,\beta)$ with $\beta<1$ to bias
        training toward tail quantiles.
        Achieves CRPS~183.04, coverage~0.947---nearly meeting the
        0.95 target from Stage~1 alone.
  \item \textbf{Stage~2} (Exog + Adaptive sampling):
        resumes from Stage~1 weights
        and re-weights quantile sampling by the softmax of calibration
        gap magnitudes $p(\tau) \propto \exp(\gamma|g(\tau)|)$.
        Achieves the \textbf{best multi-seed result}:
        $\mathrm{CRPS} = 172.21 \pm 1.5$, coverage~0.907,
        an 18.5\% improvement over the published baseline.
  \item \textbf{Stage~3} (Exog + Adaptive + Series): adding series
        embeddings slightly degrades to CRPS~172.75 and coverage~0.895,
        confirming Stage~2 as the optimal stopping point.
\end{itemize}

\begin{table}[H]
\centering
\caption{Curriculum learning stage results. Best configuration in bold.}
\label{tab:curriculum}
\begin{tabular}{llcc}
\toprule
\textbf{Stage} & \textbf{Configuration} & \textbf{CRPS} &
\textbf{Cov-0.95} \\
\midrule
\multicolumn{4}{l}{\textit{Curriculum v1 -- feature-staging (3 seeds)}} \\
Stage 1 & Exog (baseline)           & 173.0              & ---             \\
Stage 2 & Exog + Series             & 181.0              & ---             \\
Stage 3 & Exog + Series + Adaptive  & 190.0              & 0.885           \\
\midrule
\multicolumn{4}{l}{\textit{Curriculum v2 -- sampling-first (5 seeds)}} \\
Stage 1 & Exog + Beta sampling      & $183.04\pm1.9$     & $0.947\pm0.012$ \\
\textbf{Stage 2} & \textbf{Exog + Adaptive (best)} &
  $\mathbf{172.21\pm1.5}$ & $\mathbf{0.907\pm0.008}$ \\
Stage 3 & Exog + Adaptive + Series  & $172.75\pm1.8$     & $0.895\pm0.011$ \\
\bottomrule
\end{tabular}
\end{table}

The core insight of v2 is that Beta sampling in Stage~1 establishes
well-calibrated tail behaviour as a stable starting point---the very
condition that plain adaptive sampling lacked---and that adaptive
fine-tuning in Stage~2 then compresses the distribution toward better
CRPS without sacrificing coverage.
The 0.907 residual gap motivates CQR post-processing
(Section~\ref{sec:cqr_results}).

% ---------------------------------------------------------------------------
\subsection{Novel Architectural Contributions}
\label{subsec:novel_results}

All three novel output heads are evaluated on the standard NBEATSAQCAT
backbone \emph{without} exogenous features, using the published
hyperparameters to isolate the architectural effect.

\subsubsection{Quantile-Conditioned Basis Coefficients (QCBC)}

QCBC modulates each N-BEATS block's basis coefficients by the quantile
level $\tau$, injecting quantile awareness at the representation level
rather than only at the output head.
The near-baseline CRPS ($213.22$ vs.\ $211.22$) suggests that the standard
basis-expansion output already provides sufficient quantile representational
capacity at 74.1~M parameters.
The modulation scale of $0.001$ was chosen to prevent training instability;
larger values led to divergence in preliminary experiments.

\subsubsection{Temporal Coherence Regularisation (TCR)}

TCR penalises large temporal differences between adjacent quantile forecast
steps.  The higher CRPS ($217.44$) and largest standard deviation ($\pm5.42$)
reflect instability in the adaptive penalty weight.
Fixed-weight configurations achieved CRPS~214--216, still above the
baseline.  The fundamental tension is that temporal smoothness and CRPS
sharpness are competing objectives for electricity demand, which exhibits
sharp morning and evening load ramps.

\subsubsection{Distributional Basis Expansion (DBE)}

DBE replaces the linear forecast head with a 3-component Laplace mixture
model.  DBE achieves $\mathrm{CRPS} = 212.86$ with a best-seed result
of~210.99---marginally below the paper baseline.
The blend weight $\alpha$ converged to $\approx 0.08$ across seeds,
indicating a modest corrective signal from the mixture head.
The 937~K additional parameters add limited value on the 35-country
EMHIRES dataset, which may lack sufficient distributional heterogeneity
for the mixture to learn distinct modes.

\begin{table}[H]
\centering
\caption{Novel architectural head results (5~seeds unless noted).
  Published baseline CRPS~$= 211.22$.}
\label{tab:novel}
\begin{tabular}{llccc}
\toprule
\textbf{Contribution} & \textbf{Architecture} & \textbf{CRPS} &
\textbf{Cov-0.95} & \textbf{Params} \\
\midrule
Baseline (paper)  & NBEATSAQCAT      & 211.22             & ---             & 74.1 M \\
QCBC              & NBEATSAQQCBC     & $213.22\pm2.65$    & $0.927\pm0.008$ & 74.1 M \\
TCR               & NBEATSAQCAT+TCR  & $217.44\pm5.42$    & $0.923\pm0.016$ & 74.1 M \\
DBE               & NBEATSAQCAT+DBE  & $212.86\pm1.50$    & $0.924\pm0.015$ & 75.1 M \\
\bottomrule
\end{tabular}
\end{table}

% ---------------------------------------------------------------------------
\subsection{Alternative Loss Functions and Non-Crossing Constraints}
\label{subsec:alt_loss_results}

Table~\ref{tab:alt_loss} shows that all alternative loss functions
underperform the MQN baseline.
The smooth Huber variant replaces the sharp kink of the pinball loss
with a smooth quadratic region near
$y = \hat{q}_\tau$, hypothesised to improve gradient flow; the $4.3\%$
CRPS degradation suggests the kink serves a useful role in precisely
localising the quantile boundary.
The non-crossing cumsum architecture forces strict quantile monotonicity
by construction but reduces flexibility, confirming the finding
of~\cite{chernozhukov2010} that post-hoc sorting is preferable to hard
architectural constraints for CRPS optimisation.

\begin{table}[H]
\centering
\caption{Alternative loss function and constraint results
  (single seed unless noted).}
\label{tab:alt_loss}
\begin{tabular}{llcc}
\toprule
\textbf{Method} & \textbf{Backbone} & \textbf{CRPS} & \textbf{Cov-0.95} \\
\midrule
MQN (baseline)          & NBEATSAQCAT  & 211.22 & ---   \\
Smooth Huber pinball    & NBEATSAQCAT  & 220.34 & 0.898 \\
Smooth Arctan pinball   & NBEATSAQFILM & 249.91 & 0.916 \\
Hybrid CRPS-Pinball     & NBEATSAQFILM & 250.73 & 0.913 \\
Non-crossing (cumsum)   & NBEATSAQCAT  & 217.06 & 0.922 \\
\bottomrule
\end{tabular}
\end{table}

% ---------------------------------------------------------------------------
\subsection{Conformalized Quantile Regression Calibration}
\label{sec:cqr_results}

\noindent
\textbf{Theoretical background.}
CQR~\cite{romano2019cqr} provides distribution-free coverage guarantees
under exchangeability.
For lower/upper quantile predictions $\hat{q}_{\alpha/2}(x_i)$ and
$\hat{q}_{1-\alpha/2}(x_i)$, the nonconformity score on calibration
data is (Eq.~\ref{eq:cqr_score_method}):
\[
s_i = \max\!\left\{
  \hat{q}_{\alpha/2}(x_i) - y_i,\;
  y_i - \hat{q}_{1-\alpha/2}(x_i)
\right\}.
\]
Let $\hat{s}_{1-\alpha}$ be the empirical $(1-\alpha)$-quantile of
$\{s_i\}_{i=1}^{n}$.  The calibrated prediction interval
(Eq.~\ref{eq:cqr_interval_method}) becomes:
\[
\mathcal{C}(x)=\Bigl[
  \hat{q}_{\alpha/2}(x)-\hat{s}_{1-\alpha},\;
  \hat{q}_{1-\alpha/2}(x)+\hat{s}_{1-\alpha}
\Bigr].
\]
This correction is model-agnostic, requires no retraining, and guarantees
$\geq 1-\alpha$ marginal coverage on exchangeable test data.

\medskip
\noindent
\textbf{Experimental setup.}
CQR is applied to the Stage~2-v2 model using a deliberate seasonal
calibration/test split to maximise the distribution-shift challenge:
\begin{itemize}
  \item \textbf{Calibration set}: January--June~2018
        (winter and spring demand patterns).
  \item \textbf{Test set}: July--December~2018
        (summer and autumn---systematically different distribution).
\end{itemize}

\begin{table}[H]
\centering
\caption{CQR calibration results. Calibration: Jan--Jun 2018.
  Test: Jul--Dec 2018. Target coverage: 0.95.}
\label{tab:cqr}
\rowcolors{2}{gray!8}{white}
\begin{tabular}{lccccc}
\toprule
\textbf{Method} & \textbf{CRPS$\downarrow$} &
\textbf{Cov-0.05$\uparrow$} & \textbf{Cov-0.50} &
\textbf{Cov-0.95$\uparrow$} & \textbf{Interval width} \\
\midrule
Raw (no calibration)   & 126.631 & 0.0538 & 0.4816 & 0.9389 & ---           \\
\rowcolor{green!10}
CQR (per-quantile)     & 126.663 & 0.0398 & 0.4955 & \textbf{0.9593}
                        & $+69.74$~MW \\
AQ-NBEATS (paper)      & 211.22  & ---    & ---    & ---    & ---           \\
\bottomrule
\end{tabular}
\end{table}

\medskip
\noindent
\textbf{Interpretation.}
CQR achieves $95.93\%$ empirical coverage at $\alpha=0.95$ despite the
winter-to-summer distribution shift, confirming its robustness to seasonal
non-stationarity.  The slight over-coverage ($0.9593 > 0.95$) is the
conservative, theoretically expected behaviour of conformal methods
(they guarantee $\geq 1-\alpha$, not exact equality).
The CRPS change from $126.631$ to $126.663$ is negligible ($+0.03\%$),
indicating that the $+69.74$~MW average interval expansion is a modest
price for a \emph{formal} coverage guarantee.

The symmetric conformal method (constant expansion of both bounds) achieves
only coverage~0.914, demonstrating that the asymmetric per-quantile CQR
score is essential under seasonal-shift conditions.

\begin{figure}[H]
\centering
\includegraphics[width=\textwidth]{cqr_results.png}
\caption{CQR calibration analysis for Stage~2-v2 (Exog + Adaptive).
  \textbf{Top-left}: Calibration curve (empirical vs.\ nominal coverage);
  CQR (green) lies significantly closer to the ideal diagonal than the
  raw model (red), confirming improved marginal calibration.
  \textbf{Top-right}: Per-quantile coverage error
  (empirical $-$ nominal); CQR reduces the systematic under-coverage of
  the raw model at lower quantiles (Q0.01--Q0.40) to within $\pm1\%$.
  \textbf{Bottom}: Empirical vs.\ nominal coverage bar chart across the
  full quantile set; CQR bars align closely with the nominal target.
  \textbf{Table inset}: Summary metrics confirming coverage
  and near-unchanged CRPS.}
\label{fig:cqr_results}
\end{figure}

% ===========================================================================
\section{Evaluation and Discussion}
\label{sec:eval}
% ===========================================================================

\subsection{Analysis of Model Evaluation Curve}
\label{subsec:eval_curves}

\subsubsection{Validation CRPS and Coverage Convergence}

\begin{figure}[H]
\centering
\includegraphics[width=0.95\textwidth]{crps_evolution_figure_final.png}
\caption{Validation CRPS evolution across 15 training epochs for the
  best Stage~2-v2 configuration (Exog + Adaptive sampling).
  The dashed vertical line marks the selected best checkpoint (epoch~4).
  Shaded bands show $\pm 1$ standard deviation across 5 seeds.}
\label{fig:crps_evolution}
\end{figure}

Figure~\ref{fig:crps_evolution} reveals a characteristic two-phase
convergence pattern present across all configurations:
\begin{enumerate}
  \item \textbf{Rapid improvement (epochs 1--5):} The model learns
        dominant daily and weekly demand seasonalities.  Exogenous
        variants converge $\approx$30\% faster because direct access to
        temperature eliminates the need to infer weather-driven demand
        from lagged history.
  \item \textbf{Slow refinement (epochs 6--15):} Tail quantiles and
        country-specific features are progressively refined, with smaller
        but consistent CRPS improvements.
\end{enumerate}

The best checkpoint at epoch~4 (minimum validation CRPS) demonstrates
that early stopping provides meaningful regularisation against
over-fitting to the training distribution.

\subsubsection{Epoch Duration Curve}

\begin{figure}[H]
\centering
\includegraphics[width=0.95\textwidth]{epoch_duration_curve.png}
\caption{Training epoch duration across 15 epochs for the Stage~2-v2 configuration. The curve shows stable computational performance with slight increase in later epochs due to adaptive sampling overhead.}
\label{fig:epoch_duration}
\end{figure}

Figure~\ref{fig:epoch_duration} demonstrates consistent training efficiency across epochs, with average epoch duration of 18.2 minutes. The adaptive sampling mechanism introduces minimal computational overhead (~5\% increase) while providing significant performance benefits.

\subsubsection{Learning Rate Schedule and Warm-Up}

The inverse square-root schedule with a 400-batch warm-up was critical for
the exogenous models.
Runs \emph{without} warm-up diverged in 3 out of 5 seeds due to instability
in the embedding lookup gradients during the first mini-batches.
The warm-up phase linearly increases the learning rate from $10^{-6}$ to
$10^{-4}$, allowing the embedding table and the exogenous encoder to
stabilise before the full gradient magnitude is applied.

\subsubsection{Gradient Stability Across Architectural Variants}

The TCR and DBE variants exhibit higher gradient variance during the first
3 epochs, explaining their larger CRPS standard deviations across seeds.
The QCBC small-initialisation strategy (modulation scale $= 0.001$)
successfully prevents early instability while allowing quantile modulation
to develop gradually.

% ---------------------------------------------------------------------------
\subsection{Model Performance Comparison}
\label{subsec:perf_comparison}

\subsubsection{Internal Ranking of All Configurations}

Table~\ref{tab:internal} ranks all evaluated configurations by CRPS.
The curriculum~v2 Stage~2 model (CRPS $= 172.21$) is the best
reproducible multi-seed result; the single-seed exogenous run
(CRPS $= 168.46$) represents the best observed but unconfirmed result.

\begin{table}[H]
\centering
\caption{Consolidated internal ranking by CRPS. Best reproducible
  result in bold.}
\label{tab:internal}
\begin{tabular}{clccc}
\toprule
\textbf{Rank} & \textbf{Configuration} & \textbf{CRPS} &
\textbf{Cov-0.95} & \textbf{$\Delta$CRPS} \\
\midrule
1  & Stage~2-v2: Exog+Adaptive        & $\mathbf{172.21}\pm1.5$ & 0.907 & $-18.5\%$   \\
2  & Stage~3-v2: Exog+Adaptive+Series & $172.75\pm1.8$          & 0.895 & $-18.2\%$   \\
3  & Exog only (Exp~20, 1 seed)       & $168.46^*$              & 0.919 & $-20.3\%^*$ \\
4  & Exog+Series (5 seeds)            & $174.09\pm0.95$         & 0.912 & $-17.6\%$   \\
5  & Series embedding only            & $206.51\pm0.88$         & 0.936 & $-2.2\%$    \\
6  & AQ-NBEATS paper baseline         & 211.22                  & ---   & ---         \\
7  & DBE novel head                   & $212.86\pm1.50$         & 0.924 & $+0.8\%$    \\
8  & QCBC novel head                  & $213.22\pm2.65$         & 0.927 & $+0.9\%$    \\
9  & Non-crossing (cumsum)            & 217.06                  & 0.922 & $+2.8\%$    \\
10 & Smooth Huber loss                & 220.34                  & 0.898 & $+4.3\%$    \\
\bottomrule
\multicolumn{5}{l}{\footnotesize $^*$Single seed; multi-seed mean may differ.}
\end{tabular}
\end{table}

\subsubsection{External Comparison with Published Benchmarks}

Table~\ref{tab:external} places our best results in the context of the
full published benchmark from~\cite{smyl2025anyquantile}.
Our best multi-seed model (CRPS $= 172.21$) surpasses AQ-ESRNN
(CRPS $= 195.94$) by approximately $12\%$.
This is noteworthy because AQ-ESRNN uses a hybrid architecture
(exponential smoothing + dilated RNN) specifically designed for
electricity demand, while our enhancement operates on the simpler
N-BEATS backbone.  The result empirically confirms the hypothesis
of~\cite{smyl2025anyquantile} that exogenous features are more impactful
than ESRNN architectural advantages on this benchmark.

\begin{table}[H]
\centering
\caption{External comparison with published benchmarks
  \protect\cite{smyl2025anyquantile}.
  Our results are highlighted in the lower section.}
\label{tab:external}
\rowcolors{2}{gray!8}{white}
\begin{tabular}{lcccc}
\toprule
\textbf{Model} & \textbf{CRPS} & \textbf{N-CRPS} & \textbf{MAPE (\%)} &
\textbf{vs.\ AQNBEATSPlusPlus} \\
\midrule
WaveNet              & 293.38 & 2.52 & 3.39 & $+38.9\%$ \\
ES                   & 325.80 & 2.64 & 3.48 & $+54.3\%$ \\
AQ-ESRNN (paper)     & 195.94 & 1.72 & 2.32 & $-7.2\%$  \\
AQNBEATSPlusPlus (paper)    & 211.22 & 1.84 & 2.47 & ---       \\
\midrule
\rowcolor{green!12}
Exog only, Exp~20 (1 seed)        & 168.46 & --- & 2.11 & $-20.3\%^*$ \\
\rowcolor{green!12}
Stage~2-v2: best multi-seed       & 172.21 & --- & 2.13 & $-18.5\%$   \\
\rowcolor{green!12}
Stage~2-v2 + CQR (post-hoc)      & 172.21 & --- & ---  & (cov.~0.960)\\
\bottomrule
\multicolumn{5}{l}{\footnotesize $^*$Single seed.}
\end{tabular}
\end{table}

% ---------------------------------------------------------------------------
\subsection{Qualitative Analysis of Probabilistic Forecasts}
\label{subsec:qualitative}

\begin{figure}[H]
\centering
\includegraphics[width=0.95\textwidth]{cqr_effect_figure.png}
\caption{Qualitative forecast comparison for a representative winter
  period (Jan~2018).
  \textbf{Grey area}: 90\% prediction interval (Q5--Q95 band).
  \textbf{Green line}: median forecast (Q50).
  \textbf{Blue line}: observed demand.
  The CQR-calibrated model (right panel) produces a well-centred interval
  reliably containing the observed demand, reducing both under-coverage
  events and excessive over-coverage width visible in the raw model
  (left panel).}
\label{fig:qualitative}
\end{figure}

\textbf{Country-level observations.}
Exogenous features provide the most visible improvements for high
temperature-sensitivity countries: Germany, France, and Norway.
For these, winter cold spells drive demand spikes that purely autoregressive
models predict too conservatively (wide intervals, poorly centred).
The exogenous-augmented model narrows the PI while better centring it on
the observed peak.
For low temperature-sensitivity countries (Iceland, Luxembourg) the
improvement is smaller, consistent with their variability coefficients
$v_y < 5\%$.

\textbf{Seasonal coverage balance.}
Baseline models tend to under-cover in winter (strongest temperature
effects) and over-cover in summer (more regular patterns).
The curriculum~v2 model exhibits a more balanced seasonal coverage profile:
beta-sampling in Stage~1 improves winter tail coverage, and adaptive
sampling in Stage~2 reduces summer over-coverage.

\textbf{Quantile crossing.}
Post-hoc sorting reduces CRPS by approximately 2--5\% for AQNBEATSPlusPlus-based
models, consistent with the smaller benefit reported in
\cite{smyl2025anyquantile} relative to AQ-ESRNN, where the sigmoid
non-linearity caused 30\% crossing degradation.

% ---------------------------------------------------------------------------
\subsection{Limitations of the Best Model}
\label{subsec:limitations}

\textbf{1.~Coverage gap.}
Stage~2-v2 achieves coverage~0.907 against the 0.95 nominal target
before CQR.  This gap of $\approx$4.3 percentage points is closed to
0.960 by CQR post-processing (Section~\ref{sec:cqr_results}).

\textbf{2.~Single-seed exogenous result.}
Exp~20 (CRPS $= 168.46$) derives from a single training run.
Given the $\pm1.5$ standard deviation observed in multi-seed exogenous
experiments, the true multi-seed mean is likely in the range 166--171.
This result should therefore be interpreted as a lower-bound estimate
of achievable multi-seed performance with full exogenous features,
pending additional production runs.

\textbf{3.~Reanalysis vs.\ operational weather data.}
Our experiments use historical ERA5 reanalysis as a proxy for operational
weather forecast inputs.  In deployment, temperature forecasts carry
lead-time-dependent uncertainty.  The true operational CRPS improvement
may be smaller, particularly at the 36--48~hour horizon where
forecast errors are larger.

\textbf{4.~Holiday and event modelling.}
The model has no explicit holiday indicator, relying on calendar embeddings
to partially capture holiday effects.  Residual holiday-period errors
represent the main remaining source of tail miscalibration, visible at
the Q0.01 and Q0.99 extremes in the per-quantile coverage error plot
(Figure~\ref{fig:cqr_results}, top-right).

% ===========================================================================
\section{Graphical Interface}
\label{sec:gui}
% ===========================================================================

The model is deployed as a production-ready web application exposing the
Stage~2-v2 checkpoint and CQR calibrator through a FastAPI back-end and a
React~18 / TypeScript front-end.
The interface targets energy grid operators and analysts who need to
generate, inspect, and export probabilistic demand forecasts without direct
access to the Python codebase.
The application implements JWT authentication, per-user API keys, a
RESTful forecast endpoint (\texttt{POST /api/forecast/forecast}), and four
primary views described below.

\begin{figure}[H]
\centering
\includegraphics[width=0.92\textwidth]{ui_login.png}
\caption{Login page. Users authenticate with email and password; the
  JWT token is stored client-side and attached to all subsequent API
  requests as a \texttt{Bearer} header.  Credentials are validated
  against the PostgreSQL user table.  The form displays inline
  validation feedback for malformed inputs and an incorrect-credentials
  error on failed authentication.}
\label{fig:ui_login}
\end{figure}

\begin{figure}[H]
\centering
\includegraphics[width=0.92\textwidth]{ui_dashboard.png}
\caption{Home / Dashboard page. Displays a summary of the most recently
  generated forecast: four KPI cards (peak demand, average demand,
  data quality score, number of loaded variables), a 48-hour forecast
  chart with Q5--Q95 prediction bands and Q50 median line, a grid-capacity
  alert when any quantile forecast exceeds the configurable threshold
  (default 11\,000~MW), and a per-variable data quality panel populated
  from the uploaded CSV.}
\label{fig:ui_dashboard}
\end{figure}

\begin{figure}[H]
\centering
\includegraphics[width=0.92\textwidth]{ui_forecast.png}
\caption{Forecast page. Left panel: configuration controls including
  variable selector (auto-populated from the uploaded CSV), custom
  quantile slider (1--99\%), forecast horizon selector (24H / 48H / 72H),
  additional quantile overlay checkboxes (Q5, Q25, Q50, Q75, Q90), and
  forecast start date-time picker.  Right panel: Recharts area chart
  showing all selected quantile lines and the capacity reference line,
  together with an Export CSV button that downloads the full quantile
  forecast table.}
\label{fig:ui_forecast}
\end{figure}

\begin{figure}[H]
\centering
\includegraphics[width=0.92\textwidth]{ui_data.png}
\caption{Data management page. Drag-and-drop CSV upload accepting the
  \texttt{ds, unique\_id, y} long-format schema.
  Upon upload the page parses all variables, computes per-variable
  statistics (min, max, mean, standard deviation, missing count, quality
  score), and displays them in a sortable table with colour-coded quality
  bars.  Clicking the eye icon opens a modal with a mini time-series
  chart and a raw data preview for the selected variable.
  A Download Template button exports a correctly formatted sample CSV.}
\label{fig:ui_data}
\end{figure}

